% Part: methods
% Chapter: proofs
% Section: using-definitions

\documentclass[../../../include/open-logic-section]{subfiles}

\begin{document}

\olfileid{mod}{prf}{def}

\olsection{నిర్వచనాలను ఉపయోగించడం}

నిరూపణలో అవసరమయ్యే అన్ని నిర్వచనాలు మీకు తెలిసి ఉండాలని, వాటిని సరిగ్గా వర్తింపజేయగలగాలని ముందే చెప్పాం. ఇది చాలా ముఖ్యమైన విషయం; అందుకే కొంచెం వివరంగా చూద్దాం. లక్షణాలు, సంబంధాలను సంక్షిప్తంగా చెప్పడానికి నిర్వచనాలు ఉపయోగపడతాయి. ప్రవేశపెట్టే సంక్షిప్త పదాన్ని \emph{నిర్వచ్యపదం (definiendum)} అంటారు; అది దేనిని సంక్షిప్తం చేస్తుందో దాన్ని \emph{నిర్వచక భాగం (definiens)} అంటారు. నిరూపణలలో నిర్వచ్యపదాన్ని ఎలా ప్రవేశపెట్టామో తరచుగా మళ్లీ చూడాల్సి వస్తుంది. ఎందుకంటే నిరూపణను కొనసాగించడానికి నిర్వచక భాగం యొక్క తార్కిక నిర్మాణాన్ని (నిర్వచించిన పదం సూచించే పూర్తి రూపాన్ని) ఉపయోగించాలి. నిర్వచనాలను విప్పి చూస్తే తార్కిక పని జరిగే అసలు భాగం స్పష్టమవుతుంది.

ఒక ఉదాహరణతో మొదలుపెడదాం. మీరు ఈ క్రింది విషయాన్ని నిరూపించాలనుకుందాం:

\begin{prop}
ఏ సమితులు $A$, $B$లకైనా $A \cup B = B \cup A$.
\end{prop}

నిరూపణ మొదలుపెట్టాలంటేనే రెండు సమితులు ఒకటే అనడం అంటే ఏమిటో తెలుసుకోవాలి; అంటే ఈ సమీకరణలోని ``$=$''కు సమితుల సందర్భంలో అర్థం ఏమిటో తెలుసుకోవాలి. రెండు సమితులకు ఒకే \tetoken{మూలకాలు}{!!{element}s} ఉంటేనే అవి ఒకటేనని నిర్వచనం. కాబట్టి మనం విప్పాల్సిన నిర్వచనం ఇది:

\begin{defn}
$A$, $B$ సమితులు \emph{ఒకటే}, $A = B$, అప్పుడూ అప్పుడే $A$లోని ప్రతి \tetoken{మూలకం}{!!{element}}, $B$లో \tetoken{ఒక మూలకం}{!!a{element}} అయి, వ్యతిరేక దిశలోనూ అలాగే ఉంటుంది.
\end{defn}

ఈ నిర్వచనంలో $A$, $B$లను ఏవైనా సమితులకు స్థానభర్తీలుగా వాడుతున్నాం. ఇది నిర్వచించేది---\emph{నిర్వచ్యపదం}---``$A = B$'' అనే వ్యక్తీకరణ; $A = B$ సత్యం అయ్యే షరతును చెప్పడం ద్వారా నిర్వచిస్తుంది. ఆ షరతు---``$A$లోని ప్రతి \tetoken{మూలకం}{!!{element}}, $B$లో \tetoken{ఒక మూలకం}{!!a{element}} అయి, వ్యతిరేక దిశలోనూ అలాగే ఉంటుంది''---\emph{నిర్వచక భాగం}.\footnote{ఈ ప్రత్యేక సందర్భంలో---కొంచెం గందరగోళంగా ఉన్నా!---$A = B$ అయితే, ఎడమ, కుడి వైపుల వేర్వేరు అక్షరాలు వాడినప్పటికీ $A$, $B$ ఒకే సమితి. అయితే ఆ సమితిని నిర్దేశించే మార్గాలు వేర్వేరు కావచ్చు; అందుకే నిర్వచనం నిస్సారమైనది కాదు.} ఆ షరతు నెరవేరితే, అప్పుడే, నిర్వచనం ప్రకారం $A = B$ సత్యం (ఈ ``అప్పుడూ అప్పుడే''ను సంక్షిప్తంగా ``iff'' అంటాం).

నిర్వచనాన్ని వర్తింపజేసేటప్పుడు అందులోని $A$, $B$లను మీ ప్రస్తుత సందర్భానికి సరిపోల్చాలి. మన సందర్భంలో $A \cup B = B \cup A$ సత్యం కావాలంటే, ప్రతి $z \in A \cup
B$ కూడా $B \cup A$లో ఉండాలి; వ్యతిరేక దిశలోనూ అలాగే. ప్రతిపాదనలోని $A \cup
B$ వ్యక్తీకరణ నిర్వచనంలోని~$A$ పాత్రను, $B
\cup A$ నిర్వచనంలోని~$B$ పాత్రను పోషిస్తాయి. $A$, $B$లు నిర్వచనంలోనూ, నిరూపిస్తున్న ప్రతిపాదనలోనూ వేర్వేరు పాత్రలలో వాడబడ్డాయి; వాటిని కలపకుండా జాగ్రత్తపడాలి. ఉదాహరణకు $A$లోని ప్రతి \tetoken{మూలకం}{!!{element}}, $B$లో \tetoken{ఒక మూలకం}{!!a{element}} అని, అలాగే వ్యతిరేక దిశలోనూ చూపితే ప్రతిపాదన నిరూపితమైందనుకోవడం తప్పు---అది $A = B$ను చూపుతుంది గానీ $A \cup B = B \cup A$ను కాదు. (పైగా $A$, $B$ ఏ సమితులైనా కావచ్చు; వాటి గురించి ఏమీ అనుకోకపోతే $A$, $B$ వేర్వేరు సమితులై ఉండవచ్చు, కనుక ఆ మార్గం ఉపయోగపడదు.)

ఈ నిరూపణలో సమ్మేళనం వంటి సమితి-సిద్ధాంత భావనలు ఉన్నాయి; కాబట్టి నిరూపణ ఎలా సాగాలో అర్థం కావాలంటే $\cup$ సంకేతం అర్థం కూడా తెలుసుకోవాలి. కొన్నిసార్లు ఒక నిర్వచనాన్ని విప్పితే మరిన్ని నిర్వచనాలను విప్పాల్సి వస్తుంది. ఉదాహరణకు $A \cup B$ నిర్వచనం $\Setabs{z}{z \in A
  \text{ or } z \in B}$. కాబట్టి $x \in A \cup
B$ నిరూపించాలంటే, $\cup$ నిర్వచనం ప్రకారం $x \in \Setabs{z}{z \in A \text{ or } z \in B}$ నిరూపించాలి. తరువాత $x \in \Setabs{z}{\dots z\dots}$ అప్పుడూ అప్పుడే $\dots
x\dots$ అని కూడా గుర్తుంచుకోవాలి. అందుచేత $\Setabs{z}{\dots z \dots}$ సంకేతాన్ని ఇంకా విప్పితే చూపాల్సింది $x \in
A$ లేదా $x \in B$. కాబట్టి ``$A \cup B$లోని ప్రతి \tetoken{మూలకం}{!!{element}}, $B \cup A$లో కూడా \tetoken{ఒక మూలకం}{!!a{element}}'' అన్నది నిజానికి ``ప్రతి $x$కు, $x \in
A$ లేదా $x \in B$ అయితే, $x \in B$ లేదా $x \in A$'' అని అర్థం. ప్రతిపాదనలోని నిర్వచనాలన్నిటినీ పూర్తిగా విప్పితే చూపాల్సింది ఇది:

\begin{prop}
ఏ సమితులు $A$, $B$లకైనా: (a) ప్రతి $x$కు, $x \in A$ లేదా $x \in
B$ అయితే $x \in B$ లేదా $x \in A$; అలాగే (b) ప్రతి $x$కు, $x \in B$ లేదా $x \in A$ అయితే $x \in A$ లేదా $x \in B$.
\end{prop}

ముఖ్యమైనది ఏమిటంటే నిరూపణను నిర్మించడంలో నిర్వచనాలను విప్పడం తప్పనిసరి భాగం. సరిగ్గా చేయడం కొన్నిసార్లు కష్టం: నిర్వచనంలోని చరరాశులను, మీరు నిరూపిస్తున్న వాదనలోని పదాలను వేరుగా గుర్తించి సరిగ్గా సరిపోల్చాలి. విజయం సాధించాలంటే ప్రశ్న ఏమి అడుగుతోందో, అందులోని అన్ని పదాల అర్థం ఏమిటో తెలుసుకోవాలి---తరచుగా ఒక్కటికంటే ఎక్కువ నిర్వచనాలను విప్పాల్సి ఉంటుంది. దిగువన ఉన్నవంటి సరళ నిరూపణలలో సమాధానం నిర్వచనాల నుంచే దాదాపు వెంటనే వస్తుంది. అయితే ఎప్పుడూ ఇంత సులభంగా ఉండదు.

\begin{prob}
$A \cap B \neq \emptyset$ అని నిరూపించమని అడిగారని అనుకుందాం. ఇక్కడ వచ్చే అన్ని నిర్వచనాలను విప్పండి; అంటే ``$\cap$'', ``='', ``$\emptyset$'' ప్రస్తావించకుండా ఈ ప్రకటనను మరో రూపంలో చెప్పండి.
\end{prob}

\end{document}
